\chapter{Schistosoma Antibody Test}

\section*{What Is This Test?}
Schistosoma antibody testing detects antibodies to blood flukes that cause schistosomiasis. It is useful when infection is suspected after freshwater exposure in an endemic region, particularly when egg detection is negative or cannot be performed reliably.

\section*{Alternative Names}
\begin{itemize}
  \item Schistosomiasis Serology
  \item Bilharzia Antibody Test
  \item Schistosome IgG
  \item Schistosoma Antibodies
\end{itemize}
\section*{How It Is Performed}
Serum is tested by an enzyme immunoassay, indirect immunofluorescence, or another validated serologic method. The assay may detect antibodies to one or more Schistosoma species and generally does not distinguish active from previous infection.

\section*{Why It Is Ordered}
\begin{itemize}
  \item Evaluation of eosinophilia, hematuria, abdominal symptoms, or hepatosplenic disease after residence or travel in an endemic area
  \item Screening people with relevant freshwater exposure before immunosuppression or other high-risk treatment
  \item Supporting diagnosis when stool or urine microscopy is negative despite compatible exposure and symptoms
\end{itemize}

\section*{Interpretation and Limitations}
A positive result supports previous exposure or infection but does not reliably prove active infection or treatment failure. Antibodies may persist after cure, and early infection or immunosuppression can produce a negative result. Stool and urine microscopy, antigen or molecular testing, eosinophil count, and organ-specific evaluation may be needed. Results should be interpreted with travel history and exposure details.
\section*{Specimen and Preparation}
The test uses serum and generally requires no fasting. Provide countries visited, freshwater exposure dates, residence history, prior schistosomiasis treatment, eosinophil results, and immune-suppressing medicines. Different assays may detect different species and may not distinguish current from previous infection.

\section*{Risks and Safety}
Physical risk is limited to venipuncture. A positive antibody result should not automatically trigger repeated treatment because antibodies may persist after cure. Severe abdominal swelling, jaundice, blood in urine or stool, seizures, or neurologic symptoms require clinical evaluation for organ complications.

\section*{Understanding Results and Follow-up}
A positive result supports exposure or infection but does not reliably prove active infection or treatment failure. A negative result may occur early or with immune suppression. Follow-up may include stool or urine microscopy, antigen or molecular testing, eosinophil count, urinalysis, and organ-specific assessment, interpreted with exposure history and symptoms.

\section*{Regulatory Status and Authorization Date}
This chapter describes a test category, not a specific commercial product. FDA, CDSCO, EU IVDR, and MHRA approval or registration dates therefore cannot be assigned without the assay manufacturer, product name, instrument, and intended use. For laboratory-developed tests, an individual agency approval date may not exist. See the Preface for the interpretation of regulatory status.

\clearpage
